% Avsnitt 8.7, ur lectures/L08/README.md.

\section{Sammanfattning}
\label{c8:sec:review}

\needspace{6\baselineskip}
\subsection*{Det här ska du kunna nu}
\begin{itemize}
  \item Förklara vad en ändlig tillståndsmaskin är, och peka ut dess tillstånd, övergångar,
    ingångar och utgångar i en given krets.
  \item Konstruera en Mooremaskin för hand, från specifikation via tillståndsdiagram och
    tillståndstabell till Karnaughhärledda ekvationer, och bygga resultatet i CircuitVerse.
  \item Modellera samma maskin i VHDL med en uppräknad typ \code{state_t} och en
    \code{case}-process.
  \item Känna igen en Mealymaskin på sekunden och säga vad dess utgång kostar och köper, en
    klockcykels latens mot ett extra tillstånd, och varför Moore är standardvalet här och i
    CAN-projektet.
  \item Komponera en tillståndsmaskin med tidigare byggda subkomponenter i stället för att skriva
    om dem. I capstonen betyder det tre tidigare kapitels moduler runt en maskin som är den enda
    nya logiken, med skift- och räknaridiom från ytterligare ett.
\end{itemize}

\needspace{6\baselineskip}
\subsection*{Frågor att testa dig själv med}
\begin{enumerate}
  \item Vad gör en krets till en tillståndsmaskin, snarare än bara ett register med lite logik
    omkring?
  \item Vad är \code{X} exakt i den handkonstruerade maskinen, och varför måste det vara en puls på
    en klockcykel snarare än knappens nivå?
  \item Varför är den maskinens tillståndskodning en Graykod? Vad kan gå fel med en vanlig
    binärräkning?
  \item En räknare och ett skiftregister är tekniskt sett också tillståndsmaskiner. Vad skiljer dem
    från maskinen som konstruerades här?
  \item Vilka signaler får en Mooreutgång bero på, och vilka får en Mealyutgång bero på? Hur avgör
    du, givet en enda rad VHDL som driver en utgång, vilken av dem du tittar på?
  \item Varför behöver Mooreversionen i \secref{c8:sec:mealy}:s jämförelse ett extra tillstånd, och
    varför ligger dess utgång exakt en klockcykel efter Mealyversionens?
  \item Varför läser \code{LED_PROCESS} i \file{fsm_led.vhd} bara \code{state}, aldrig
    \code{button_edge_s2} eller \code{button_n} direkt? Vad skulle ändras om den gjorde det?
  \item Varför går inte \code{din} i \code{seq_detect_mealy} genom en synkroniserare, till skillnad
    från \code{reset_n}?
  \item Vad köper \code{when others} dig i ett \code{case} för tillståndsövergångar, när
    \code{state_t} bara har de värden du deklarerat?
  \item Varför går timern i capstone-mottagaren sexton gånger så fort som baudhastigheten snarare
    än en gång per bit, och varför gör det att designen tål en sändare vars klocka skiljer sig något
    från din, att sampla mitt i en bit?
  \item Capstone-sändaren använder en helbitstimer och behöver ingenting mitt i biten. Vad har den
    som mottagaren inte har, som gör skillnaden?
\end{enumerate}

\needspace{6\baselineskip}
\subsection*{Blicken framåt}
Det här är sista kapitlet i bokens första del. Över åtta kapitel har vi gått från grindar och
Karnaughdiagram, via register, metastabilitet och själva VHDL-språket, till räknare, skiftregister
och timers, och nu ändliga tillståndsmaskiner, i varje steg med resonemang om kretsen för hand innan
den skrevs i syntetiserbar VHDL och kördes på en riktig FPGA. Varje mönster som byggts längs vägen,
dubbelvippsynkroniseraren, kodstilen med \code{process} och \code{case}, och nu Mooremaskinen, har
återanvänts i det här kapitlets egna exempel, och följer med oförändrat in i det som kommer härnäst.

Härnäst är \chapref{c9:ch}, den första praktiska tentamen: tre timmar, individuell, vid datorn med
GHDL, över allt de åtta första kapitlen har byggt. Efter den börjar grupprojektet, där en
CAN-kontroller konstrueras i VHDL över tio kapitel och avslutas mot en riktig mikrokontroller över
SPI.
